\documentclass[a4paper,10pt]{article}
\usepackage[utf8]{inputenc}
\usepackage{graphicx}
\usepackage{subfigure}
\usepackage{fancyhdr}
\usepackage{listings}
\usepackage{ gensymb }


%opening

\pagestyle{fancy} %eigener Seitenstil
\fancyhf{} %alle Kopf- und Fußzeilenfelder bereinigen
\fancyhead[L]{SEP Abgabegespräch SS2015/16} %Kopfzeile links
\fancyhead[C]{} %zentrierte Kopfzeile
\fancyhead[R]{\includegraphics[width=1cm]{../../template/tugraz_bw}} %Kopfzeile rechts
\renewcommand{\headrulewidth}{0.4pt} %obere Trennlinie
\fancyfoot[C]{} %Seitennummer
\renewcommand{\footrulewidth}{0.4pt} %untere Trennlinie


\lstset{language=C++, basicstyle=\ttfamily\tiny}  

\begin{document}

\section*{Monster-Verwaltung}

Schreibe ein einfaches Programm um Monster für ein Computerspiel zu verwalten. Die Monster können mit verschiedenen Namen und Health Points (HP) erzeugt werden. Außerdem kann allen Monstern Schaden zugefügt werden.

\begin{description}
 \item[spawn $<$monsterName$>$ $<$monsterHP$>$] Erstellt ein neues Monster mit gegebenem Namen und Health points.
 \item[list] Listet alle momentan lebenden Monster mit Namen und HP auf. Das Format der Ausgabe ist ''$<$monsterName$>$ : $<$monsterHP$>$''
 \item[damage $<$damageValue$>$] Fügt allen Monstern, die lebendig sind $damageValue$ Schadenspunkte zu. Falls dadurch die Health Points auf 0 oder darunter ($\leq 0$) sinken sollte, wird ''$<$monsterName$>$ is removed!'' ausgegeben.
\end{description}

Bei Fehlerfällen (Falsche Parameter, kein Monster existiert...) soll nichts ausgegeben werden.
Bei der Eingabe von \texttt{quit} oder EOF soll das Programm beendet werden.

\subsection*{Beispiel}
\begin{verbatim}
> spawn Ozma 65535
> spawn Goomba 1
> list
Ozma:65535
Goomba:1
> damage 1
Goomba is removed!
> list
Ozma:65534
> damage 99
> list
Ozma:65435
> damage 60000
> list
Ozma:5435
> damage 6000
Ozma is removed!
> list
> quit

\end{verbatim}

\subsection*{Aufruf}
\texttt{./monster-management}

\subsection*{Bewertung}
Dokumentation \& Programmierstil,  Dokumentation \& Programmierstil, und Struktur \& Korrektheit (je 1 Punkt)
\newpage


\end{document}
